\chapter{Background}
\label{chap:background}

% TODO: Establish notation for curves, moduli spaces, and stacks.
% TODO: Add the needed intersection theory, Hodge classes, Gromov--Witten
% theory, virtual fundamental classes, and localization, with citations.

\section{Category Theory}
\label{section:categoryThry}

\begin{Def}
    A \term{category} \(\cat{C}\) consists of 
    \begin{itemize}
        \item A class \emph{Obj} of entities called objects.
        \item A class \emph{Arw} of entities called arrows.
        \item Assignments
        \[\textit{source}:\textit{Arw}\to\textit{Obj},\word{and}\textit{target}:\textit{Arw}\to\textit{Obj}.\]
        These assignments are represented in the obvious way: \(A\xrightarrow{f}B\) indicates that \(f\) is an arrow with \emph{source} \(A\) and \emph{target} \(B\). 
        \item An assignment 
        \[
        1:\textit{Obj}\to\textit{Arw}
        \]
        which, given an object \(A\in\cat{C}\), yields an arrow \(1_A\) satisfying \(A\xrightarrow{1_A}A\). In other words, the \emph{source} and \emph{target} assignments of the distinguished arrow \(1_A\) are \(A\) itself.
        \item A partial composition
        \[
        \textit{Arw}\x\textit{Arw}\to\textit{Arw}.
        \]
        Given two arrows \(A\xrightarrow{f}B_1\) and \(B_2\xrightarrow{g}C\), they are composable when \(B_1\) and \(B_2\) are the same object. The resulting arrow has the form \(A\xrightarrow{fg}C\).
    \end{itemize}
    Additionally the category \(\cat{C}\) must satisfy the following conditions:
    \begin{itemize}
        \item Suppose we are given a diagram as follows:
        \[
        A\xrightarrow{f}B\xrightarrow{g}C\xrightarrow{h}D.
        \]
        We require that \((fg)h=f(gh)\). In other words, composition is associative.
        \item Consider an arrow \(f\) and the two identity arrows as follows:
        \[
        A\xrightarrow{1_A}A\xrightarrow{f}B\xrightarrow{1_B}B.
        \]
        We require that \(f1_B=f=1_Af\).
    \end{itemize}
\end{Def}

\begin{Rmk}
We say ``class'' and ``assignment'' to allow for the possibility that \emph{Obj} and \emph{Arw} are not sets. If they were, then an assignment would be the same as a function. We also denote, \(\Mor_{\cat{C}}(A,B)\) for the class of all arrows with source \(A\) and target \(B\). 
\end{Rmk}

\subsection{Examples of categories}

\begin{Ex}
    In these categories, every object has extra properties and arrows are functions which respect this property. Rather than making the notion of \emph{concrete category} precise, we will list some examples:
    \begin{table*}[h]
        \centering
        \begin{tabular}{lll}\toprule
            Category & Objects & Arrows\\ \midrule
            \(\cat{Set}\)& sets & functions\\
            \(\cat{Grp}\)& groups & morphisms\\
            \(\cat{AbGrp}\)&  Abelian groups & morphisms\\
            \(\cat{Rng}\)&  rings & morphisms\\
            \(\cat{Top}\)&  topological spaces & continuous functions\\
            \(\cat{Vect}_{k}\)&  vector spaces over a field \(k\) & linear maps\\
       \bottomrule
        \end{tabular}
        \end{table*}
\end{Ex}

\begin{Def}
    For a category \(\cat{C}\), \(\cat{D}\) is a \term{subcategory} of \(\cat{C}\) if the class of objects (resp. arrows) of \(\cat{D}\) is a subset of the class of objects (resp. arrows) of \(\cat{C}\).\par
    Moreover \(\cat{D}\) is a \term{full subcategory} of \(\cat{C}\) if it's a subcategory and for all objects \(X,Y\) of \(\cat{D}\), 
    \[
    \Mor_{\cat{D}}(X,Y)=\Mor_{\cat{C}}(X,Y).
    \]
\end{Def}

An examples of an important category that is studied in algebraic geometry is the following:

\begin{Ex}
    
\end{Ex}


